\section{Discussion and Limitations}

The mature experiment answers two questions. First, seam supervision is not merely an early-training artifact: the component under-count improvement persists after matched extended optimization. Second, more training alone does not eliminate the trade-off. The three-class formulation gives a sparse seam class equal status in the objective, but it does not distinguish a reliable inter-instance gap from an uncertain low-contrast boundary. The worse HD95 indicates that a small number of distant errors dominate the surface tail.

The image-level CMR should be interpreted cautiously. It detects component under-counting, but it does not identify which anatomical pair merged, and fragmentation can reduce CMR while worsening component-count MAE. The diagnostic is therefore paired with overlap, recall, boundary, surface, and count metrics. Statistical uncertainty and repeated-seed variation are not yet reported. The dataset also comes from a single pediatric hand-radiograph task, and the age distribution is imbalanced. These limitations prevent a claim of general anatomical topology preservation.

The current implementation uses a three-class head and one backbone. A stronger model-agnostic version should decouple the prior from the architecture: a confidence-weighted seam map would act on ordinary binary logits, and a two-sided support loss would penalize indiscriminate erosion. Age and sex should modulate confidence rather than define a fixed spatial template because hand size, pose, and ossification patterns are not globally aligned. Claims of plug-and-play transfer require matched experiments on at least U-Net, a multiscale decoder, and nnU-Net.

Finally, the seam target is derived from training annotations. The inference path receives only the radiograph, but the distinction must remain explicit to avoid ground-truth leakage. The held-out clean-test-v2 split remains locked until losses, thresholds, and checkpoint rules are frozen.
